\subsection{Will the fire leave the second brush?}
We tested a region selected by the original stroke, where a single yes or no could inform an edit. Around the second brush of radius 0.12 we use the immediately adjacent shell, out to radius 0.26, one flame interaction radius beyond the brush. Does any currently cold site in this shell ignite within the window? The center follows the second stroke in each take. The rule, all cuts, and all horizons were fixed before regional scoring. The same two trajectories were reused. A native evaluator can stop on the first regional ignition.

All 18 queries were nonempty. Ten contained an ignition and eight did not. For this query the weighted rule certified 8,816 individual site--window pairs. Unit weights certified 6,537. Every window still contained an undecided site, including all eight with no ignition. At frame 240 of the first take, for example, the weighted rule certified 1,842 of 1,863 sites through the next scan. None ignited. The 21 unresolved sites nevertheless prevented an early answer. All 54 region-specific ID sets exactly matched the restrictions of their full-domain counterparts. Region-specific evaluation and native fallback together cost a median paired $1.40\times$ the ordinary first-event-or-end replay, starting with the checkpoint loaded; including file decoding increased this ratio to $3.05\times$. The native stepping component had a median of 17.7 ms. These additive costs use an optimistic native first-event cutoff from the trace and exclude an integrated interface.

We then tested two cheap ways to certify a positive answer on the same regions: enough residual heat to ignite at the next scan, and a lower bound on future heating from guaranteed surviving graph donors. The regional outcomes were already known when this follow-up was designed. Both rules abstained on every window, including the ten positives. A surviving donor prefix was available in all twelve multistep windows. This result bounds what these particular witnesses resolve; an ignition not proved by them cannot be assigned a causal source on that basis.

\balance
\subsection{Where the ordinary replay can stop}
We instrumented a fresh pair of continuous native takes to timestamp completion of the ignition scan within a step. Earlier steps must complete to provide the next state. The terminal step's graph, fuel and emission tail is needed only if the state must continue. Within the same execution, the median time for complete label queries was 71.8 ms through the full terminal step and 64.9 ms through its scan (median paired ratio 0.896). For the regional question, stopping at the first in-shell event or the horizon end, the corresponding medians were 18.6 and 11.85 ms (paired ratio 0.637). All recorded frame events agreed with the uninstrumented-by-this-timer reference. The timer measures prefixes of ordinary computation. A callable early-return implementation would also pay for dispatch, materializing the result and testing regional membership. These single-host paired comparisons use a separate run from the 17.7 ms regional timing above.
